% =============================================================================
\section{Transformadas de Laplace y Z}
% =============================================================================

\begin{frame}{Transformada de Laplace}
  \begin{definicion}[transformada bilateral]
    \[
      X(s) = \TL{x(t)} = \int_{-\infty}^{\infty} x(t)\, \e^{-st} \dd t,
      \qquad s = \sigma + \j\omega \in \ROC
    \]
  \end{definicion}
  \begin{ejemplo}
    \vspace{-0.5\baselineskip}
    \begin{align*}
      \e^{-at}\escalon{t} &\laplace \frac{1}{s + a}, & \Re\{s\} &> -a \\
      -\e^{-at}\escalon{-t} &\laplace \frac{1}{s + a}, & \Re\{s\} &< -a
    \end{align*}
  \end{ejemplo}
  \begin{importante}
    \small Misma expresión, señales distintas: la \alert{región de convergencia} es parte de la transformada.
  \end{importante}
\end{frame}

\begin{frame}{Polos, ceros y estabilidad}
  \begin{columns}[c]
    \begin{column}{0.5\textwidth}
      \[
        H(s) = \frac{s + 1}{(s + 2)(s^2 + 2s + 5)}
      \]
      \begin{itemize}
        \item Cero: $s = -1$
        \item Polos: $s = -2$, \ $s = -1 \pm \j 2$
        \item<2-> Sistema causal: $\ROC = \{\Re\{s\} > -1\}$
        \item<3-> La ROC contiene el eje $\j\omega$ $\Rightarrow$ \alert{BIBO estable}
      \end{itemize}
    \end{column}
    \begin{column}{0.46\textwidth}
      \centering
      \begin{tikzpicture}
        \begin{axis}[plano s, width=\linewidth, xmin=-3, xmax=1.5, ymin=-3, ymax=3,
            xtick={-2,-1}, ytick={-2,2}]
          \only<2->{\fill[celeste, opacity=0.25] (axis cs:-1,-3) rectangle (axis cs:1.5,3);
                    \draw[celeste, thick, dashed] (axis cs:-1,-3) -- (axis cs:-1,3);
                    \node[celeste!60!black, font=\small] at (axis cs:0.75,-2.4) {ROC};}
          \addplot[cero] coordinates {(-1,0)};
          \addplot[polo] coordinates {(-2,0) (-1,2) (-1,-2)};
        \end{axis}
      \end{tikzpicture}
    \end{column}
  \end{columns}
\end{frame}

\begin{frame}{Sistemas realimentados}
  \centering
  \begin{tikzpicture}[node distance=1cm and 1.4cm]
    \node (x) {$X(s)$};
    \node[sumador, right=of x] (s) {};
    \node[bloque, right=of s] (g) {$G(s)$};
    \node[rama, right=1.2cm of g] (r) {};
    \node[right=of r] (y) {$Y(s)$};
    \node[bloque, below=0.8cm of g] (h) {$H(s)$};
    \draw[flecha] (x) -- node[pos=0.85, above, senial] {$+$} (s);
    \draw[flecha] (s) -- node[above, senial] {$E(s)$} (g);
    \draw[thick] (g) -- (r);
    \draw[flecha] (r) -- (y);
    \draw[flecha] (r) |- (h);
    \draw[flecha] (h) -| node[pos=0.9, left, senial] {$-$} (s);
  \end{tikzpicture}

  \vspace{1em}
  \pause
  \begin{columns}[T]
    \begin{column}{0.44\textwidth}
      \begin{align*}
        E(s) &= X(s) - H(s)\,Y(s) \\
        Y(s) &= G(s)\,E(s)
      \end{align*}
    \end{column}
    \begin{column}{0.44\textwidth}
      \[
        \Rightarrow\quad \resaltar<3->{T(s) = \frac{Y(s)}{X(s)} = \frac{G(s)}{1 + G(s)\,H(s)}}
      \]
    \end{column}
  \end{columns}
\end{frame}

\begin{frame}{Respuesta en frecuencia: diagrama de Bode}
  \framesubtitle{Sistema de segundo orden
    $H(s) = \dfrac{\omega_n^2}{s^2 + 2\zeta\omega_n s + \omega_n^2}$, con $\omega_n = 1$}
  \centering
  \begin{tikzpicture}
    \begin{semilogxaxis}[bode, name=mag, width=0.48\textwidth, height=4.6cm,
        domain=0.1:10, ylabel={$\abs{H(\j\omega)}$ [\si{dB}]}, ymin=-40, ymax=15,
        legend pos=south west]
      \foreach \z in {0.1, 0.3, 0.7} {
        \addplot+[] {-10*log10((1 - x^2)^2 + (2*\z*x)^2)};
      }
      \legend{$\zeta = 0{,}1$, $\zeta = 0{,}3$, $\zeta = 0{,}7$}
    \end{semilogxaxis}
    \begin{semilogxaxis}[bode, at={(mag.east)}, anchor=west, xshift=1.6cm,
        width=0.48\textwidth, height=4.6cm,
        domain=0.1:10, ylabel={$\angle H(\j\omega)$ [\si{\degree}]},
        ymin=-180, ymax=0, ytick={-180,-135,-90,-45,0}]
      \foreach \z in {0.1, 0.3, 0.7} {
        \addplot+[] {-atan2(2*\z*x, 1 - x^2)};
      }
    \end{semilogxaxis}
  \end{tikzpicture}
\end{frame}

\begin{frame}{Transformada Z}
  \begin{columns}[c]
    \begin{column}{0.54\textwidth}
      \begin{definicion}
        \[
          X(z) = \TZ{x[n]} = \sum_{n=-\infty}^{\infty} x[n]\, z^{-n}
        \]
      \end{definicion}
      \[
        a^n \escalon{n} \ztransf \frac{1}{1 - a z^{-1}}, \quad \abs{z} > \abs{a}
      \]
      \begin{itemize}
        \item Polos en $z = 0{,}5$ y $z = 0{,}8\,\e^{\pm\j\pi/4}$
        \item Causal y estable $\Leftrightarrow$ todos los polos dentro del círculo unitario
      \end{itemize}
    \end{column}
    \begin{column}{0.42\textwidth}
      \centering
      \begin{tikzpicture}
        \begin{axis}[plano z, width=\linewidth, xmin=-1.4, xmax=1.4, ymin=-1.4, ymax=1.4,
            xtick={-1,1}, ytick={-1,1}]
          \fill[celeste, opacity=0.25, even odd rule]
            (axis cs:-1.4,-1.4) rectangle (axis cs:1.4,1.4)
            (axis cs:0,0) circle[radius=0.8];
          \addplot[cero] coordinates {(0,0)};
          \addplot[polo] coordinates {(0.5,0) (0.5657,0.5657) (0.5657,-0.5657)};
          \node[celeste!60!black, font=\small] at (axis cs:-1.05,1.15) {ROC};
        \end{axis}
      \end{tikzpicture}
    \end{column}
  \end{columns}
\end{frame}
